\begin{abstract}
Video LLMs are increasingly asked not only \emph{what} happened but \emph{when}: a grounded question such as
``how long did I use the paper cutter?'' needs an answer and the time interval of its evidence. To keep long videos
tractable, the visual memory of these models is pruned to a small fraction of its tokens, and published pruners lose
little multiple-choice accuracy at 10\% of the tokens. We show that temporal grounding needs something different:
\emph{pruning for QA keeps what is salient, but grounding needs what is everywhere}. A video LLM reads time from
timestamps that survive only together with visual tokens of their frame, so pruners that select by saliency,
uniqueness or text similarity empty frames, delete their timestamps, and fall \emph{below random token selection} on
grounding. When the answer lists observed moments before the interval, the gaps are filled by a video-free prior:
asked with no video at all, the model still starts its intervals at the beginning of the window. We propose
\method, a training-free method that selects frames by the question (BOLT) and adds two parts. \emph{Coverage} gives every frame an equal token quota, chosen per
layer by the actual question's attention, so that every time anchor survives. \emph{Prior-contrastive decoding}
(PCD) writes each time value against the model's own no-video prediction. The two parts behave differently across
answer formats, which we study directly. Under the standard one-line grounding prompt, Coverage alone matches the
uncompressed model with a quarter of the visual tokens on a general and a grounding-tuned model of the same family
(QVHighlights, 61.7 \vs 60.5 and 67.3 \vs 67.1 mIoU) and beats random selection (\key{$+3.5$} and \key{$+2.0$}),
while PCD has no effect; on the full Charades-STA and QVHighlights sets of TimeLens-Bench, Coverage at 31\% of the
tokens is lossless on both models and beats both random selection and a lower-resolution input of the same budget. In reasoning-style grounded QA (S-EMBER), PCD is what matters: with 10\% of the tokens
\method improves on the uncompressed model by \key{$+11.2$} and \key{$+10.7$} mIoU under the same prompt, and
matches an oracle that sees only the gold evidence on grounded answers. Asking
for a timeline of observations instead of a single interval costs single-moment grounding about 30 mIoU, so the
answer format is a first-order variable in evaluating compressed memories. We evaluate every method against
video-free floors and report where the method helps less: on short, densely sampled videos (Charades-STA,
NExT-GQA) the gain of Coverage is small at moderate budgets and grows as the budget shrinks (\key{$+6.8$} accuracy
over random selection at 5\% on NExT-GQA), and PCD helps only where the answer format lets the prior choose the times.
\end{abstract}
